\documentclass[greek, histinit, english, twosided]{book}
\usepackage{microtype}
\usepackage{lmodern}
\raggedbottom
\setlength{\emergencystretch}{2em}
\usepackage{used-packages}
\usepackage{packages/ntua-thesis-titlepages}
\usepackage{amsmath}
\usepackage{makecell}
\usepackage{listings}
\usepackage{enumitem}
\usepackage{tabularx}
\usepackage{placeins}
\usepackage{parcolumns}
\usepackage{float}
\usepackage{wrapfig}
\usepackage{booktabs}
\lstdefinestyle{mystyle}{
    keywordstyle=\color{magenta},
    numberstyle=\tiny\color{codegray},
    stringstyle=\color{codepurple},
    basicstyle=\ttfamily\footnotesize,                
    captionpos=b,
    upquote=true
}
\usepackage[nottoc]{tocbibind}
\usepackage{graphicx}
\usepackage{listings}
\usepackage{xcolor}
\definecolor{codegreen}{rgb}{0,0.6,0}
\definecolor{codegray}{rgb}{0.5,0.5,0.5}
\definecolor{codepurple}{rgb}{0.58,0,0.82}
\definecolor{backcolour}{rgb}{0.95,0.95,0.92}
\lstset{
    language=C,
    backgroundcolor=\color{backcolour},   
    commentstyle=\color{codegreen},
    keywordstyle=\color{magenta},
    numberstyle=\tiny\color{codegray},
    stringstyle=\color{codepurple},
    basicstyle=\ttfamily\footnotesize,
    breakatwhitespace=false,         
    breaklines=true,                 
    captionpos=b,                    
    keepspaces=true,                 
    numbers=left,                    
    numbersep=5pt,                  
    showspaces=false,                
    showstringspaces=false,
    showtabs=false,                  
    tabsize=2
}
\usepackage{url}
\usepackage{hyperref}
\usepackage[backend=biber,style=ieee]{biblatex}
\addbibresource{related_work.bib}
\fancypagestyle{plain}{
  \fancyhf{}
  \fancyfoot[C]{\thepage}
}

\begin{document}

\pagenumbering{arabic}
\setcounter{page}{1}
\pagestyle{empty}
\maketitle

\fancyhf{}
\fancyfoot[C]{\thepage}
\renewcommand{\headrulewidth}{0.4pt}
\renewcommand{\footrulewidth}{0pt}
\fancyhead[LO]{}
\fancyhead[RE]{}
\fancyhead[LE]{\nouppercase\leftmark}
\fancyhead[RO]{\nouppercase\rightmark}

\cleardoublepage
\pagestyle{fancy}
\chapter*{Περίληψη}
\thispagestyle{fancy}
\addcontentsline{toc}{chapter}{Περίληψη}

Ο πολλαπλασιασμός πίνακα με διάνυσμα (General Matrix--Vector Multiplication, GEMV) αποτελεί τον βασικό υπολογιστικό πυρήνα της εκτέλεσης (inference) βαθυών νευρωνικών δικτύων, καθώς για κάθε γραμμή εξόδου απαιτείται η μεταφορά και η επεξεργασία μεγάλου όγκου βαρών. Η δομημένη αραιότητα N:M, όπου κάθε ομάδα M βαρών περιέχει ακριβώς N μη μηδενικά στοιχεία, μειώνει τον απαραίτητο υπολογιστικό όγκο βαρών διατηρώντας ένα κανονικό μοτίβο που μπορεί να αξιοποιηθεί αποδοτικά σε υλικό. Η παρούσα διπλωματική εργασία παρουσιάζει τη σχεδίαση, την υλοποίηση και την αξιολόγηση ενός επιταχυντή Sparse GEMV 2:M για την πλατφόρμα (FPGA) AMD Alveo U280, ο οποίος υποστηρίζει τους λόγους αραιότητας 2:4, 2:8, 2:16 και 2:32 και έχει τη δυνατότητα να αλλάζει λόγο κατά τη λειτουργία μέσω ενός κωδικού 2 bit που μεταφέρεται μέσα στη ροή δεδομένων και δηλώνει τον λόγο που εφαρμόζεται στα βάρη. Στην προτεινόμενη αρχιτεκτονική ο αριθμός των πολλαπλασιαστών είναι σταθερός και μεταβάλλεται μόνο η διάρκεια συσσώρευσης κάθε γραμμής: 64 ανεξάρτητοι συσσωρευτές, οργανωμένοι σε 8 πυρήνες με 8 υποπυρήνες ο καθένας, επεξεργάζονται τα μη μηδενικά στοιχεία ανά δύο πάνω σε ένα σταθερό παράθυρο ενεργοποιήσεων (activations), ώστε τα DSPs να παραμένουν απασχολημένα σε κάθε λόγο αραιότητας χωρίς δέντρο άθροισης. Ένας ενταμιευτής επανάληψης (replay buffer) φορτώνει το διάνυσμα ενεργοποιήσεων (activations vector) μία φορά και το επαναχρησιμοποιεί για όλες τις γραμμές. Ο επιταχυντής ενσωματώνεται ως RTL kernel στο Vitis, με HLS data movers πάνω σε 17 ψευδοκανάλια High Bandwidth Memory (HBM), και συγκρίνεται με μια πυκνή (dense) σχεδίαση αναφοράς που χρησιμοποιεί τους ίδιους movers, την ίδια μνήμη και το ίδιο ρολόι. Και οι δύο σχεδιάσεις παράγουν αποτελέσματα ακριβή σε επίπεδο bit στην πραγματική κάρτα U280, ενώ η αραιή σχεδίαση αλλάζει λόγο αραιότητας στη μέση ενός υπολογισμού χωρίς επαναφόρτωση δεδομένων. Στα 300~MHz, η αραιή σχεδίαση υπολογίζει το ίδιο GEMV 2.00$\times$, 3.96$\times$, 7.93$\times$ και 15.81$\times$ ταχύτερα από την πυκνή για τους λόγους 2:4, 2:8, 2:16 και 2:32 αντίστοιχα και μειώνει την ενέργεια που χρειάζεται η κάρτα για τον υπολογισμό ίδιων διαστάσεων εισόδου κατά 1.77$\times$ έως 13.99$\times$, ενώ η επιτάχυνση διατηρείται και στα 325~MHz (2.02$\times$, 4.01$\times$, 8.00$\times$ και 15.92$\times$ αντίστοιχα). Το κόστος της συγκεκριμένης υλοποιημένης αρχιτεκτονικής είναι 41\% περισσότερα LUTs και 8192 επιπλέον πολυπλέκτες F7, ενώ η αραιή σχεδίαση ικανοποιεί τους χρονικούς περιορισμούς έως τα 325~MHz, έναντι 350~MHz για την πυκνή. Οι μετρήσεις δείχνουν επίσης ότι το 84--86\% της κατανάλωσης ενέργειας της κάρτας είναι στατικό, επομένως η εξοικονόμηση προέρχεται από την ταχύτερη ολοκλήρωση του υπολογισμού και όχι από χαμηλότερη κατανάλωση ισχύος. Για τη μελέτη της κλιμάκωσης εξετάζονται, σε υλοποίηση out-of-context, όλοι οι συνδυασμοί πυρήνων και υποπυρήνων με γινόμενο 64, οι οποίοι έχουν τις ίδιες διεπαφές και τον ίδιο αριθμό πράξεων ανά κύκλο ρολογιού: η διάταξη $8\times8$ δίνει την υψηλότερη συχνότητα, ενώ η αύξηση του αριθμού των πυρήνων κοστίζει ταυτόχρονα επιφάνεια και συχνότητα, με τη διάταξη $64\times1$ να απαιτεί 38\% περισσότερα LUTs και να λειτουργεί σε 15\% χαμηλότερη συχνότητα.
% TODO (όταν ολοκληρωθεί): μία πρόταση για την οικογένεια διατάξεων 4x4 έως 4x32 (128 έως 1024 DSP).
% TODO (μόνο αν το εγκρίνει ο καθηγητής): μία πρόταση για HBM έναντι DDR.
Τα αποτελέσματα δείχνουν ότι, όταν η αραιότητα χρησιμοποιείται για να μειώσει τα δεδομένα βαρών ανά γραμμή ενώ όλοι οι πολλαπλασιαστές παραμένουν απασχολημένοι, η ρυθμαπόδοση κλιμακώνεται σχεδόν γραμμικά με τον λόγο αραιότητας.

\textbf{Λέξεις-κλειδιά} --- FPGA, Επιταχυντής υλικού, Δομημένη αραιότητα N:M, GEMV, High Bandwidth Memory, Alveo U280, Αναδιάρθρωση κατά τη λειτουργία, Κλιμακωσιμότητα

\chapter*{Abstract}
\addcontentsline{toc}{chapter}{Abstract}

General matrix--vector multiplication (GEMV) is the core computational kernel of deep neural network inference, since every output row requires a large volume of weights to be transferred and processed. N:M structured sparsity, in which every group of M weights contains exactly N non-zero elements, reduces the volume of weights that must be processed while preserving a regular pattern that can be exploited efficiently in hardware. This thesis presents the design, implementation and evaluation of a 2:M sparse GEMV accelerator for the AMD Alveo U280 FPGA platform. The accelerator supports the 2:4, 2:8, 2:16 and 2:32 sparsity ratios and can switch between them at runtime through a 2-bit code that is carried within the data stream and specifies the ratio applied to the weights. In the proposed architecture the number of multipliers is fixed and only the accumulation length of each row varies: 64 independent accumulators, organized into 8 cores of 8 blocks each, process the non-zero elements two at a time over a fixed activation window, so that the DSP blocks remain busy at every sparsity ratio without an adder tree. A replay buffer loads the activation vector once and reuses it for all rows. The accelerator is integrated as an RTL kernel in Vitis, with HLS data movers over 17 High Bandwidth Memory (HBM) pseudo-channels, and is compared against a dense reference design that uses the same data movers, the same memory and the same clock. Both designs produce bit-exact results on the physical U280 card, and the sparse design changes its sparsity ratio in the middle of a computation without reloading data. At 300~MHz, the sparse design computes the same GEMV $2.00\times$, $3.96\times$, $7.93\times$ and $15.81\times$ faster than the dense design for the 2:4, 2:8, 2:16 and 2:32 ratios, respectively, and reduces the energy the card consumes for a computation of identical input dimensions by $1.77\times$ to $13.99\times$, while the speed-up is preserved at 325~MHz ($2.02\times$, $4.01\times$, $8.00\times$ and $15.92\times$, respectively). The cost of the implemented architecture is 41\% more LUTs and 8,192 additional F7 multiplexers, and the sparse design meets its timing constraints up to 325~MHz, compared with 350~MHz for the dense design. The measurements also show that 84--86\% of the card's energy consumption is static; the savings therefore come from completing the computation sooner rather than from lower power consumption. To study scalability, all combinations of cores and blocks whose product is 64 are examined in out-of-context implementation; these share the same interfaces and the same number of operations per clock cycle. The $8\times8$ configuration achieves the highest frequency, whereas increasing the number of cores costs both area and frequency, with the $64\times1$ configuration requiring 38\% more LUTs and operating at a 15\% lower frequency.
% TODO (when complete): one sentence on the 4x4 to 4x32 configuration family (128 to 1024 DSPs).
% TODO (only if approved by the supervisor): one sentence on HBM versus DDR.
The results show that when sparsity is used to reduce the weight data per row while all multipliers remain busy, throughput scales almost linearly with the sparsity ratio.

\textbf{Keywords} --- FPGA, Hardware accelerator, N:M structured sparsity, GEMV, High Bandwidth Memory, Alveo U280, Runtime reconfiguration, Scalability

\chapter*{Ευχαριστίες}
\addcontentsline{toc}{chapter}{Ευχαριστίες}

Θα ήθελα να ευχαριστήσω ιδιαίτερα τον επιβλέποντα καθηγητή μου, κ. Δημήτριο Σούντρη, για την καθοδήγηση και την υποστήριξή του καθ' όλη τη διάρκεια εκπόνησης της διπλωματικής μου εργασίας στο Εργαστήριο Μικροϋπολογιστών και Ψηφιακών Συστημάτων.\\
Ακόμη, θα ήθελα να ευχαριστήσω τον υποψήφιο διδάκτορα Ηλία Παπαλάμπρου και τον διδάκτορα Δημοσθένη Μασούρο για την άψογη συνεργασία και τις πολύτιμες συμβουλές τους.\\
Ευχαριστώ με όλη μου την καρδιά την οικογένειά μου: τη μητέρα μου, τον πατέρα μου και τα δύο μου αδέλφια, για την αμέριστη στήριξη και την αγάπη τους, ειδικά κατά τη διάρκεια των σπουδών μου.\\
Τέλος, ευχαριστώ ολόψυχα τους φίλους μου, ιδιαίτερα τον Αντρέα, τον Ματθαίο, τον Αχιλλέα, τον Χριστόδουλο και τον Αιμίλιο, για τη συμπαράσταση και την υπομονή τους σε κάθε βήμα αυτής της κοινής διαδρομής.

\begin{flushright}
Σώζος Κούλας\\
Οκτώβριος 2026
\end{flushright}

\tableofcontents
\clearpage

\listoffigures

% \chapter*{Εκτεταμένη Ελληνική Περίληψη}
\markboth{Εκτεταμένη Ελληνική Περίληψη}{Εκτεταμένη Ελληνική Περίληψη}
\addcontentsline{toc}{chapter}{Εκτεταμένη Ελληνική Περίληψη}

\section*{Εισαγωγή}

Η ραγδαία ανάπτυξη των υπολογιστικών υποδομών cloud και η αυξανόμενη χρήση FPGAs ως επιταχυντών υπολογισμού έχουν εισάγει νέες και ουσιαστικές προκλήσεις ασφάλειας. Στο παραδοσιακό μοντέλο απειλής, οι επιθέσεις side-channel απαιτούν φυσική πρόσβαση στη συσκευή-στόχο, εξειδικευμένο εξοπλισμό μέτρησης και άμεση ηλεκτρική σύνδεση με το κύκλωμα \cite{standaert2009introduction}. Η ενσωμάτωση των FPGAs σε περιβάλλοντα πολλαπλών χρηστών αναιρεί θεμελιωδώς αυτές τις υποθέσεις. Σε ένα κοινόχρηστο FPGA, πολλαπλοί χρήστες εκτελούν εργασίες ταυτόχρονα στο ίδιο φυσικό κύκλωμα. Ενώ η χωρική απομόνωση μεταξύ ενοικιαστών επιβάλλεται μέσω pblocks, το κοινόχρηστο δίκτυο διανομής ισχύος (Power Distribution Network, PDN) παραμένει ένα αδύνατο σημείο. Η μεταγωγική δραστηριότητα ενός χρήστη προκαλεί διακυμάνσεις τάσης που διαδίδονται μέσω του PDN και είναι ανιχνεύσιμες από οποιοδήποτε άλλο στοιχείο στο ίδιο κύκλωμα \cite{ahmed2022multi,zhao2018fpga}. Ένας κακόβουλος χρήστης μπορεί να αναπτύξει έναν αισθητήρα τάσης εντός του ολοκληρωμένου κυκλώματος και να ανακτήσει κρυπτογραφικά κλειδιά από την κατανάλωση ισχύος ενός θύματος που εκτελεί κρυπτογράφηση, χωρίς καμία φυσική παρέμβαση. Η παρούσα διπλωματική εργασία αντιμετωπίζει ακριβώς αυτή την απειλή. Προτείνεται και αξιολογείται μια ελαφριά τεχνική απόκρυψης ισχύος βασισμένη σε κυκλώματα Enhanced Ring Oscillator with Latch (ERO-L), τα οποία παράγουν θόρυβο ανεξάρτητο από τα δεδομένα στο PDN, στοχεύοντας στην υποβάθμιση του λόγου σήματος προς θόρυβο που είναι διαθέσιμος στον επιτιθέμενο. Αξιολογούνται τρεις στρατηγικές ενεργοποίησης αυξανόμενης πολυπλοκότητας, η αποτελεσματικότητα των οποίων επαληθεύεται μέσω εκτεταμένων πειραμάτων σε πραγματικό υλικό, χρησιμοποιώντας ως πρωτεύουσα μετρική την εξέλιξη κατάταξης κλειδιού.

\section*{Θεωρητικό Υπόβαθρο}

\subsection*{AES}

Η κρυπτογράφηση αποτελεί θεμελιώδη μηχανισμό προστασίας της εμπιστευτικότητας δεδομένων. Το Πρότυπο Προχωρημένης Κρυπτογράφησης (AES) είναι ο πλέον διαδεδομένος αλγόριθμος συμμετρικής κρυπτογράφησης παγκοσμίως, τυποποιημένος από το NIST το 2001 ως διάδοχος του DES \cite{fips2001advanced}. Λειτουργεί σε μπλοκ 128 bits υποστηρίζοντας κλειδιά 128, 192 και 256 bits, με 10, 12 ή 14 γύρους μετασχηματισμών αντίστοιχα \cite{daemen2002design}. Κάθε γύρος κρυπτογράφησης, εκτός του τελευταίου, αποτελείται από τέσσερις πράξεις: SubBytes (μη γραμμική αντικατάσταση bytes μέσω του πίνακα S-box), ShiftRows (κυκλική ολίσθηση γραμμών), MixColumns (γραμμικός μετασχηματισμός στηλών) και AddRoundKey (συνδυασμός κατάστασης με υπο-κλειδί μέσω XOR) \cite{fips2001advanced}. Για την παρούσα εργασία κρίσιμη είναι η ακολουθία AddRoundKey–SubBytes: το αποτέλεσμα αυτής της σύνθεσης αποτελεί το ενδιάμεσο μέγεθος που εκμεταλλεύεται η επίθεση CPA, καθώς η εξαρτημένη από τα δεδομένα κατανάλωση ισχύος κατά τον υπολογισμό του S-box συσχετίζεται στατιστικά με το μυστικό κλειδί \cite{lo2017power,brier2004correlation}.

\subsection*{Αρχιτεκτονική FPGA και Ασφάλεια}

Τα FPGAs (Field Programmable Gate Arrays) είναι επαναπρογραμματιζόμενα ψηφιακά κυκλώματα που αποτελούνται από πλέγμα λογικών μπλοκ (CLBs), καθένα εκ των οποίων περιέχει πίνακες αναζήτησης (LUTs) και flip-flops, συνδεδεμένα μέσω προγραμματιζόμενου δικτύου δρομολόγησης. Η εξέλιξη των FPGA τα έχει μετατρέψει από απλές συστοιχίες λογικής σε σύνθετα συστήματα με ενσωματωμένες μνήμες RAM, μπλοκ DSP και υποσυστήματα επεξεργαστή \cite{boutros2021fpga}. 

Σύγχρονα FPGAs σε πλατφόρμες cloud επιτρέπουν κοινή χρήση μιας συσκευής από πολλαπλούς χρήστες. Αυτή η κοινή χρήση εισάγει νέες ευπάθειες ασφάλειας: σε αντίθεση με επεξεργαστές γενικής χρήσης, τα FPGAs επιτρέπουν στους χρήστες να αναπτύσσουν αυθαίρετο υλικό, ενώ μοιράζονται φυσικούς πόρους όπως το PDN. Ένας κακόβουλος χρήστης μπορεί να εκμεταλλευτεί αυτό το κοινόχρηστο μέσο για επιθέσεις side-channel, εισαγωγή σφαλμάτων ή επιθέσεις άρνησης υπηρεσίας (denial of service) \cite{ahmed2022multi}.

\subsection*{Επιθέσεις Side-Channel και CPA}

Οι επιθέσεις side-channel (SCAs) εκμεταλλεύονται φυσικές πληροφορίες που διαρρέουν από την υλοποίηση ενός συστήματος, αντί να επιτίθενται στον ίδιο τον κρυπτογραφικό αλγόριθμο \cite{gattu2020power}. Τα κανάλια ισχύος εκμεταλλεύονται παραλλαγές κατανάλωσης ισχύος κατά τη διάρκεια υπολογισμών, τα χρονικά κανάλια παραλλαγές χρόνου εκτέλεσης και τα ηλεκτρομαγνητικά κανάλια εκπομπές από ηλεκτρονικά εξαρτήματα \cite{gandolfi2001electromagnetic}. Η ανάπτυξη FPGAs σε υποδομές ψλοθδεισάγει νέες ευπάθειες που επεκτείνουν τις SCAs πέρα από παραδοσιακές πλατφόρμες ενσωματωμένων συστημάτων \cite{zhao2018fpga}.

Η Correlation Power Analysis (CPA) εισήχθη επίσημα από τους Brier, Clavier και Olivier \cite{brier2004correlation} και βασίζεται στο μοντέλο Hamming Weight: η κατανάλωση ισχύος θεωρείται ανάλογη του αριθμού των bits που τίθενται σε 1 κατά τη διάρκεια μιας πράξης \cite{lo2017power}. Η επίθεση εκτελείται σε τέσσερις φάσεις: συλλογή ιχνών ισχύος ενώ επεξεργάζονται γνωστές είσοδοι, κατασκευή υποθέσεων Hamming Weight για κάθε πιθανό υπό-κλειδί, υπολογισμός συντελεστή συσχέτισης Pearson μεταξύ προβλεπόμενων και μετρούμενων τιμών, και ανάκτηση κλειδιού ως εκείνου με τη μέγιστη απόλυτη συσχέτιση \cite{lo2017power,brier2004correlation}.

% \begin{figure}[ht]
%     \centering
%     \includegraphics[width=1.0\textwidth]{Figures_2/attack_pipeline.png}
%     \caption*{Σύνοψη της μεθοδολογίας της CPA}
% \end{figure}

Κρίσιμο πλεονέκτημα της CPA είναι ότι επιτρέπει ανεξάρτητη επίθεση σε κάθε ένα από τα 16 bytes του AES-128, ανάγοντας το πρόβλημα αναζήτησης $2^{128}$ πιθανοτήτων σε 16 ανεξάρτητες αναζητήσεις 256 πιθανοτήτων\cite{lo2017power}. Για την αξιολόγηση αντιμέτρων χρησιμοποιείται η μετρική κατάταξης κλειδιού (key rank), που εκτιμά τη θέση του σωστού κλειδιού στη λίστα όλων των πιθανών κλειδιών ταξινομημένων κατά πιθανότητα βάσει αποτελεσμάτων CPA \cite{spielmann2023rds}. Τιμή $2^{128}$ σημαίνει μηδενική πληροφορία για τον επιτιθέμενο, τιμή 0 σημαίνει πλήρης ανάκτηση.

\section*{Μεθοδολογία}

\subsection*{Επιλογή κυκλώματος Ring-Oscillator} 
Η επιλογή κατάλληλου κυκλώματος για τη απόκρυψη ισχύος απαιτούσε δομημένη αξιολόγηση. Πέντε υποψήφια κυκλώματα ταλαντωτή δακτυλίου, ERO-cf, ERO-L-cf, ERO-L, ERO-FF-cf-L και ERO-FF-L, αξιολογήθηκαν ως προς πέντε κριτήρια: ένταση μεταγωγής, αξιοποίηση δρομολόγησης και φόρτωση, σταθερότητα ταλάντωσης, συμμόρφωση με κανόνες σύνθεσης και συμβατότητα με λειτουργικά συστήματα \cite{mahmoud2025roboost}.

Κάθε κύκλωμα υλοποιήθηκε μεμονωμένα στο FPGA και η συμπεριφορά του καταγράφηκε μέσω του αισθητήρα RDS, επιτρέποντας άμεση οπτική σύγκριση με τη βασική γραμμή AES. Τα ERO-cf, ERO-L-cf και ERO-FF-cf-L περιέχουν άμεσο συνδυαστικό βρόχο ανατροφοδότησης, προκαλώντας κρίσιμη παραβίαση DRC τύπου LUTLP-1 στο Vivado. Αυτό τα καθιστά ακατάλληλα για ανάπτυξη σε περιβάλλοντα με αυστηρή επιβολή κανόνων σύνθεσης. Μεταξύ των συμβατών επιλογών, το ERO-L υπερτερεί σαφώς έναντι του ERO-FF-L: παρουσιάζει μέση έξοδο αισθητήρα 108-109 μονάδων, έναντι ~103 μονάδων για το ERO-FF-L, λόγω χρήσης latch αντί flip-flop στο μονοπάτι ανατροφοδότησης. Ο latch λειτουργεί σε διαφανή κατάσταση, διαδίδοντας αλλαγές ταχύτερα και διατηρώντας υψηλότερη συχνότητα ταλάντωσης \cite{mahmoud2025roboost}. 

Το κύκλωμα ERO-L, αποτελείται από τέσσερα LUTs (A, B, C, D) ρυθμισμένα ως πύλες NAND και διασυνδεδεμένα σε πλήρως συνδεδεμένη τοπολογία με υψηλό fanout. Ένας latch εισάγεται στο μονοπάτι ανατροφοδότησης από τη LUT A, με εξωτερικό σήμα ενεργοποίησης που επιτρέπει άμεσο έλεγχο της ταλάντωσης. Αυτό το σχέδιο επιτυγχάνει υψηλή ένταση μεταγωγής ενώ παράλληλα αποφεύγει τον άμεσο συνδυαστικό βρόχο, εξασφαλίζοντας πλήρη συμβατότητα με τα εργαλεία σύνθεσης.

\subsection*{Στρατηγικές Ενεργοποίησης} 

Αφού επιλέχθηκε το ERO-L, το επόμενο ερώτημα ήταν πώς και πότε πρέπει να ενεργοποιείται σε σχέση με την κρυπτογράφηση AES που προστατεύει. Ο στόχος είναι να αλλοιωθούν τα ίχνη που καταγράφει ο αισθητήρας RDS του αντιπάλου. Η στρατηγική ενεργοποίησης καθορίζει τη χρονική δομή του παραγόμενου θορύβου και άρα πόσο αποτελεσματικά μπορεί ο θόρυβος να αντισταθεί στη διαδικασία συσχέτισης της CPA \cite{mangard2007power}. Αξιολογήθηκαν τρεις στρατηγικές αυξανόμενης πολυπλοκότητας:

\textbf{Στατική Ενεργοποίηση:} Το σήμα ενεργοποίησης είναι μόνιμα ενεργό καθ' όλη τη διάρκεια του πειράματος. Είναι η απλούστερη υλοποίηση, χωρίς λογική ελέγχου χρονισμού. Επειδή το ERO-L είναι ήδη ενεργό κατά τη βαθμονόμηση, η βαθμονόμηση γίνεται υπό τις πραγματικές συνθήκες λειτουργίας. Η ουσιαστική αδυναμία είναι ότι η σταθερή συνεισφορά του ERO-L αφαιρείται αυτόματα κατά τη φάση αφαίρεσης μέσης τιμής της CPA, αφήνοντας μόνο ένα αυξημένο επίπεδο θορύβου ως εμπόδιο για τον αντίπαλο \cite{mangard2007power}. Αυτό σημαίνει ότι η στατική στρατηγική δεν εισάγει ουσιαστική ίχνος-προς-ίχνος παραλλαγή, που είναι απαραίτητη για να αντισταθεί σε CPA.

\textbf{Περιοδική Ενεργοποίηση:} Το σήμα ενεργοποίησης εναλλάσσεται on/off με σταθερή συχνότητα και κύκλο λειτουργίας 50\% κατά τη διάρκεια κάθε ίχνους. Υλοποιείται μέσω αποκλειστικής λογικής ελέγχου συγχρονισμένης με το ρολόι του αισθητήρα, παραμετροποιημένης ως προς αριθμό δειγμάτων ανά ίχνος, αριθμό περιόδων και κύκλο λειτουργίας. Το σήμα ενεργοποίησης μεταφέρεται στον τομέα κύριου ρολογιού μέσω συγχρονιστή δύο flip-flop. Η στρατηγική εισάγει χρονική παραλλαγή εντός κάθε ίχνους, αλλά το μοτίβο είναι ταυτόσημο σε όλα τα ίχνη. Αυτό αποτελεί θεμελιώδη αδυναμία, καθώς ένας πληροφορημένος αντίπαλος θα μπορούσε θεωρητικά να ανακτήσει και να αφαιρέσει το περιοδικό μοτίβο. Επιπλέον, αν η επιλεγμένη περίοδος συνάδει με τη δομή γύρων AES, το αποτέλεσμα μπορεί να είναι αντίθετο από το επιθυμητό.

\textbf{Τυχαιοποιημένη Ενεργοποίηση:} Παράγεται μοναδική ψευδο-τυχαία ακολουθία ενεργοποίησης για κάθε ίχνος ανεξάρτητα. Ο αριθμός ενεργών παραθύρων, οι θέσεις τους και τα μήκη τους μεταβάλλονται από ίχνος σε ίχνος, εξασφαλίζοντας ότι η συνεισφορά ERO-L δεν έχει συνεκτική χρονική δομή. Η λογική ελέγχου, υλοποιείται ως μηχανή πεπερασμένων καταστάσεων τεσσάρων καταστάσεων (IDLE, DECIDE\_NUM\_WINDOWS, GENERATE\_WINDOWS, COLLECTING) οδηγούμενη από 16-bit LFSR μέγιστου μήκους. Ο LFSR τρέχει συνεχώς. Η τιμή που παρουσιάζει κατά την έναρξη κάθε ίχνους εξαρτάται από τον χρονισμό του σήματος data-ready του AES, παρέχοντας αποτελεσματικά διαφορετική αρχική τιμή ανά ίχνος. Επειδή το μοτίβο ενεργοποίησης είναι διαφορετικό για κάθε ίχνος, η συνεισφορά ERO-L δεν συσσωρεύεται συνεκτικά υπό τη διαδικασία συσχέτισης CPA. Αυτή είναι η θεμελιώδης ιδιότητα που καθιστά τη στρατηγική αυτή ανώτερη.

\subsection*{Πειραματική Υποδομή και Ροή Αξιολόγησης} 
Η υλοποίηση έγινε στην πλακέτα Digilent Nexys A7-100T που βασίζεται στο FPGA AMD Artix-7 XC7A100T (63400 LUTs, 126800 flip-flops, 6 CMTs με PLLs). Η αρχιτεκτονική συστήματος, εφαρμόζει φυσική απομόνωση μέσω περιορισμών Pblock: θύμα (AES-128), αντίπαλος (RDS[24], βαθμονόμηση, FIFO) και ERO-L τοποθετούνται σε διακριτές μη επικαλυπτόμενες περιοχές \cite{spielmann2023rds}. Το ERO-L βρίσκεται χωρικά μακριά και από τις δύο άλλες περιοχές, ώστε η επίδρασή του να διαδίδεται αποκλειστικά μέσω του κοινόχρηστου PDN. Λογική απομόνωση εξασφαλίζεται από την απουσία οποιωνδήποτε άμεσων συνδέσεων σήματος μεταξύ των τριών περιοχών.

\begin{figure}[ht]
    \centering
    \includegraphics[width=0.9\textwidth]{Figures_3/full_system_architecture.png}
    \caption*{Αρχιτεκτονική του συστήματος}
\end{figure}

Το σύστημα λειτουργεί σε δύο ανεξάρτητους τομείς ρολογιού: αισθητήρας στα 200 MHz και AES κυρίως στα 100 MHz, με πρόσθετα πειράματα στα 50 MHz και 20 MHz. Το ERO-L είναι ασύγχρονο ως προς αμφότερους τους τομείς και η συχνότητα ταλάντωσής του καθορίζεται από τις εσωτερικές του καθυστερήσεις διάδοσης, όχι από εξωτερικό ρολόι. Η ανάλυση CPA εκτελείται εκτός επεξεργαστή στον κεντρικό υπολογιστή με επιτάχυνση GPU, χρησιμοποιώντας τον αλγόριθμο εκτίμησης κατάταξης κλειδιού των Glowacz et al \cite{glowacz2015simpler}.


\section*{Πειραματικά Αποτελέσματα} 

\subsection*{Απροστάτευτη Υλοποίηση} 

Το πρώτο πείραμα αξιολόγησε την CPA επίθεση στο απροστάτευτο AES-128 στα 100 MHz. Η εξέλιξη του key rank, αποδεικνύει γρήγορη σύγκλιση: ξεκινώντας από τιμή 120, με 1000 ίχνη, η κατάταξη πέφτει ραγδαία μεταξύ 3.000 και 10000 ιχνών και φτάνει στο 0 περίπου στα 10000 ίχνη, σηματοδοτώντας πλήρη ανάκτηση κλειδιού. Το στενό χάσμα μεταξύ άνω και κάτω ορίου υποδηλώνει σταθερές συνθήκες μέτρησης και ισχυρή διαρροή. Αυτά τα αποτελέσματα αποτελούν τη βάση αναφοράς για όλες τις επόμενες συγκρίσεις.

\subsection*{Αποτελέσματα Στατικής Ενεργοποίησης}

Η στατική ενεργοποίηση αξιολογήθηκε με 128, 256 και 320 παράλληλα ERO-L. Η εξέλιξη του key rank αποκαλύπτει συμπεριφορά που αντίκειται στις αρχικές προσδοκίες. Τα 128 στιγμιότυπα αποτελούν την πιο αποτελεσματική στατική διαμόρφωση: η κατάταξη φτάνει περίπου στο 72 στα 10000 ίχνη χωρίς σύγκλιση. Τα 256 στιγμιότυπα παρουσιάζουν πιο σύνθετη συμπεριφορά, με ταχύτερη αρχικά αλλά βραδύτερη κατόπιν κατάβαση. Τα 320 στιγμιότυπα, η διαμόρφωση με τη μεγαλύτερη ορατή συσκότιση σε επίπεδο ίχνους, συγκλίνουν στο 0 εντός 10000 ιχνών, χειρότερα από την απροστάτευτη βασική γραμμή.

\begin{figure}[h!]
    \centering
    \includegraphics[width=0.8\textwidth]{Figures_4/keyrank_static.png}
    \caption*{Αποτελέσματα key rank για στατική ενεργοποίηση}
\end{figure}

Αυτό το αποτέλεσμα εξηγείται από τον δομικό περιορισμό της στατικής ενεργοποίησης: η σταθερή συνεισφορά ERO-L αφαιρείται από την προεπεξεργασία της CPA και η βαρύτερη φόρτωση PDN σε μεγαλύτερο αριθμό στιγμιοτύπων μετατοπίζει το σημείο λειτουργίας του αισθητήρα με τρόπους που ενδέχεται να ενισχύουν ορισμένα χαρακτηριστικά διαρροής. Το αποτέλεσμα αυτό υπογραμμίζει ότι η αύξηση αριθμού στιγμιοτύπων υπό στατική ενεργοποίηση δεν μεταφράζεται αξιόπιστα σε ισχυρότερη αντίσταση CPA.

\subsection*{Αποτελέσματα Περιοδικής Ενεργοποίησης} 

Η περιοδική ενεργοποίηση αξιολογήθηκε με τρεις διαμορφώσεις περιόδου (2, 4 και 6 πλήρεις περίοδοι ανά ίχνος 256 δειγμάτων, κύκλος λειτουργίας 50\%). Η ανάλυση εξέλιξης του key rank για τις τρεις περιόδους παρουσιάζει μη μονοτονία: η διαμόρφωση 4 περιόδων είναι η πιο επιβλαβής, καθώς όλα τα πλήθη ERO-Ls επιταχύνουν τη σύγκλιση σε σχέση με την απροστάτευτη υλοποίηση, με τα 320 να φτάνουν στο μηδέν περί τα 7000 ίχνη. Η περίοδος 64 δειγμάτων συντονίζεται δυσμενώς με τη δομή γύρων AES στα 100 MHz, με αποτέλεσμα οι μεταβάσεις on/off να εμπίπτουν με συνέπεια στις ίδιες υπολογιστικές φάσεις σε όλα τα ίχνη, ενισχύοντας τη διαρροή αντί να τη μειώνει. Πρόκειται για σαφή εφαρμογή του θεωρητικά αναμενόμενου κινδύνου.

\begin{figure}[h!]
    \centering
    \includegraphics[width=0.8\textwidth]{Figures_4/keyrank_per_2.png}
    \caption*{Αποτελέσματα key rank για περιοδική ενεργοποίηση, 128 δείγματα/περίοδο}
\end{figure}

\begin{figure}[h!]
    \centering
    \includegraphics[width=0.8\textwidth]{Figures_4/keyrank_per_4.png}
    \caption*{Αποτελέσματα key rank για περιοδική ενεργοποίηση, 64 δείγματα/περίοδο}
\end{figure}

\begin{figure}[h!]
    \centering
    \includegraphics[width=0.8\textwidth]{Figures_4/keyrank_per_6.png}
    \caption*{Αποτελέσματα key rank για περιοδική ενεργοποίηση, 42 δείγματα/περίοδο}
\end{figure}

Η διαμόρφωση 6 περιόδων είναι η πιο ευνοϊκή: τα 256 ERO-Ls διατηρούν κατάταξη άνω του 80 στα 10000 ίχνη χωρίς εμφανή τάση σύγκλισης. Η διαμόρφωση 2 περιόδων βρίσκεται ενδιάμεσα με ανάμεικτα αποτελέσματα. Το κύριο συμπέρασμα είναι ότι η σχέση μεταξύ παραμέτρων περιόδου και αμυντικής αποτελεσματικότητας δεν μπορεί να προβλεφθεί από οπτική ανάλυση ιχνών και ότι κακώς επιλεγμένη παράμετρος μπορεί να μειώσει δραστικά την ασφάλεια.

\subsection*{Αποτελέσματα Τυχαιοποιημένης Ενεργοποίησης} 
Η τυχαιοποιημένη ενεργοποίηση αξιολογήθηκε σε εκτεταμένο πείραμα 20000 ιχνών. Η εξέλιξη του key rank, επιβεβαιώνει την ισχύ της τυχαιοποιημένης ενεργοποίησης και αποκαλύπτει μια εντυπωσιακά συνεπή σχέση: όσο περισσότερα ERO-Ls, τόσο ισχυρότερη αντίσταση. Με 128 η κατάταξη φτάνει περίπου στο 25 στα 20000 ίχνη. Με 256 φτάνει περίπου στο 75 στα 20000 ίχνη χωρίς σημεία σύγκλισης. Με 320 η κατάταξη παραμένει ουσιαστικά επίπεδη κοντά στο 125 καθ' όλη τη διάρκεια του πειράματος, το ισχυρότερο αποτέλεσμα ανάμεσα σε όλες τις αξιολογούμενες στρατηγικές και αριθμούς στιγμιοτύπων.

\begin{figure}[h!]
    \centering
    \includegraphics[width=0.8\textwidth]{Figures_4/keyrank_rand.png}
    \caption*{Αποτελέσματα key rank για τυχαία ενεργοποίηση}
\end{figure}

Η βελτίωση οφείλεται στον θεμελιώδη μηχανισμό: επειδή το μοτίβο ενεργοποίησης μεταβάλλεται ανεξάρ-\\τητα για κάθε ίχνος, η συνεισφορά ERO-L δεν συσσωρεύεται συνεκτικά υπό τη διαδικασία συσχέτισης CPA ανεξαρτήτως αριθμού ERO-L. Η απουσία αποτυχιών εξαρτώμενων από ευθυγράμμιση, που κατέστησαν την περιοδική στρατηγική αναξιόπιστη, καθιστά την τυχαιοποιημένη ενεργοποίηση την πιο εύρωστη στρατηγική.

\subsection*{Επίδραση Συχνότητας Ρολογιού AES} 

Για να αξιολογηθεί η γενικευσιμότητα των αποτελεσμάτων, όλες οι στρατηγικές ενεργοποίησης δοκιμά-\\στηκαν επίσης στα 50 MHz και 20 MHz με 320 ERO-Ls. Τα αποτελέσματα αποκαλύπτουν συνεπή βελτίωση αποτελεσματικότητας σε χαμηλότερες συχνότητες για στατική και περιοδική ενεργοποίηση: η στατική ενεργοποίηση που αποτυγχάνει στα 100 MHz παρέχει ουσιαστικά πλήρη αντίσταση στα 20 MHz, και η περιοδική ενεργοποίηση ενισχύεται σε παρόμοιο βαθμό. Για τυχαιοποιημένη ενεργοποίηση η σχέση είναι λιγότερο μονότονη: καμία από τις τρεις συχνότητες δεν παρουσιάζει σύγκλιση εντός του πειραματικού προϋπολογισμού, αλλά στα 50 MHz παρατηρείται ελαφρώς βραδύτερη πτώση σε σχέση με τα 100 και 20 MHz. 

\subsection*{Συνοπτική Σύγκριση} 
Ο παρακάτω πίνακας συνοψίζει τη χρήση πόρων για τις τρεις στρατηγικές N320 στον Artix-7. Ο πίνακας ERO-L καταναλώνει 320 LUTs (0.5\%) σε όλες τις στρατηγικές. Ο ελεγκτής ενεργοποίησης μεταβάλλεται: στατική ενεργοποίηση δεν απαιτεί ελεγκτή, περιοδικός ελεγκτής προσθέτει 29 LUTs (0.05\%) και τυχαιοποιημένος ελεγκτής 129 LUTs (0.2\%). Συνολικά το σύστημα χρησιμοποιεί 3096-3225 LUTs (4.88-5.09\%), αφήνοντας άνω του 94\% των πόρων διαθέσιμο.

\begin{table}[h]
\centering
\begin{tabular}{lccc}
\hline
\textbf{Στρατηγική} & \textbf{ERO-L LUTs (\%)} & \textbf{Controller LUTs (\%)} & \textbf{Συνολικά LUTs (\%)} \\
\hline
Baseline (χωρίς ERO-L) & ---         & ---         & 2776 (4.38\%) \\
Στατική              & 320 (0.5\%) & ---         & 3096 (4.88\%) \\
Περιοδική (6-period) & 320 (0.5\%) & 29 (0.05\%) & 3125 (4.93\%) \\
Τυχαιοποιημένη          & 320 (0.5\%) & 129 (0.2\%) & 3225 (5.09\%) \\
\hline
\end{tabular}
\caption*{Απαιτούμενοι πόροι στο Artix-7 XC7A100T}
\label{tab:lut_utilization_el}
\end{table}

Ο παρακάτω πίνακας συνοψίζει την αποτελεσματικότητα της μεθόδου για N320 στα 100 MHz. Η στατική ενεργοποίηση συγκλίνει στο 0 εντός 10000 ιχνών, το ίδιο αποτέλεσμα με την απροστάτευτη βασική γραμμή, άρα αναποτελεσματική. Η περιοδική ενεργοποίηση (6 περίοδοι) διατηρεί κατάταξη ~25 στα 10000 ίχνη χωρίς σύγκλιση, ωστόσο άλλες περίοδοι είναι ενεργά επιβλαβείς. Η τυχαιοποιημένη ενεργοποίηση διατηρεί κατάταξη 120 στα 20000 ίχνη χωρίς καμία τάση σύγκλισης, η ισχυρότερη και πλέον αξιόπιστη επίδοση.

\begin{table}[h]
\centering
\begin{tabular}{lccc}
\hline
\textbf{Στρατηγική} & \textbf{Αριθμός Traces} & \textbf{Key Rank} & \textbf{Συγκλίνει?} \\
\hline
Baseline            & 10000 & 0   & Yes \\
Στατική              & 10000 & 0   & Yes \\
Περιοδική (6-period) & 10000 & 25  & No \\
Τυχαιοποιημένη          & 20000 & 120 & No \\
\hline
\end{tabular}
\caption*{Αποτελεσματικότητα των πειραμάτων}
\label{tab:obfuscation_effectiveness_el}
\end{table}

% \input{Chapters/introduction}
\chapter{Background}

\section{Deep Neural Networks (DNNs) Inference and Hardware Acceleration}
\label{sec:dnn_inference}

This section introduces the computational structure of deep neural networks (DNNs) and explains why their inference cost is mainly matrix products. It also compares the hardware platforms on which DNN inference is executed.

\subsection{Neural Networks and Inference}
\label{subsec:nn_inference}



% \begin{figure}[ht]
%     \centering
%     \includegraphics[width=0.3\textwidth]{Figures_2/aes_operations.png}
%     \caption[AES internal operations]{AES internal operations \cite{aumasson2024serious}}
%     \label{fig:aes_operations}
% \end{figure}


% \[
% H(D) = \sum_{j=0}^{m-1} d_j
% \]

% The Hamming weight power model assumes a linear relationship between the number of active bits and the instantaneous power consumption. For a device processing n-bit data word D, the power consumption W can be expressed as:
% \[
% W = a \cdot H(D) + b
% \]
% \[
% H_{\text{predicted}}(K_{\text{guess}}) = HW\bigl(S[P \oplus K_{\text{guess}}]\bigr)
% \]

% \input{Chapters/related_work}
% \input{Chapters/Llama}
% \input{Chapters/architecture_methodology}
% \input{Chapters/experimental_results}
% \input{Chapters/conclusion}


%\addcontentsline{toc}{chapter}{Bibliography}
%\bibliographystyle{unsrt} 
%\bibliography{references}
\printbibliography[heading=bibintoc]

\end{document}
